\scanpage{17}
\[
 F(\Omega)\otimes B_p/\underline f B_p\longrightarrow F(\Omega\otimes B_p/\underline f B_{\unsure{p|q}})
\]
est bijective.

\uline{Dém}, La nécessité claire, celle de (vi) résultant du \deleted{lemme 6, \unsure{du|et}} corollaire 7, et du fait qu'on peut choisir $f$ tel que $B_f$ soit CM (\unsure{donc|alors} $\underline f=(f_1\unsure{-|\ldots}f_{\unsure{n|r}})$ est nécessairement une suite \deleted{régulière} $B_p$-régulière).

\uline{Suffisance}. On \unsure{note|voit} que les conditions envisagées sont stables par passage à un $F_B$ ($B$ quotient de $A$). Cela nous permet \deleted{de} \deleted{supposer, par réc. noeth., que $F_B$ est \unsure{repr.} par $B$-module de t.f., si $B\ne A$.} \margin{(i)(ii) (iii b)(iv)}grâce à prop. \deleted{\unsure{4}}3, de nous borner à trouver un $f\in A$ non nilpotent, tel que $F_f$ soit repr.

\margin{(i)(ii) (iii a)(iv)(v)}Notons que le lemme 5 \inserted{et Remarque 5 bis} impliquent d'autre part que $F$ est \inserted{\unsure{strt|str.}} proreprésentable par un système projectif \inserted{strict} $(\unsure{P|R}_i)$, les $P_i$ de t.f., on est ramené à trouver $f$, non nilpotent, tel que le syst. $P_{if}$ soit stationnaire.

\margin{(vi)}Or cela est assuré par la \deleted{condition} condition (vi), compte tenu du cor. 11.

\begin{center}\uline{cqfd.}\end{center}

\uline{Remarque} 13 Supposons que pour tt quotient intègre $B$ de $A$, $\exists\,f\in B-0$ tel que $B_f$ soit régulier. Cette condition rend \deleted{évidemment} inutile celle que $A$ soit quotient d'un régulier, et entraîne que \unsure{dans|ds} th. 12, $(\mathrm{vi})\Longrightarrow(\mathrm v)$, \deleted{et de (vi), \unsure{=|-}}
% archive sheet number: 16
